\documentclass[10pt,aspectratio=169]{beamer}

% --- pdflatex-safe font & micro-typography ---
\usepackage[T1]{fontenc}
\usepackage[utf8]{inputenc}
\usepackage{lmodern}
\usepackage{microtype}

% --- math & symbols ---
\usepackage{amsmath,amssymb,mathtools,bm}

% --- hyperlinks (Beamer already loads hyperref; this is safe) ---
\hypersetup{hidelinks}

% --- theme ---
\usetheme[numbering=fraction]{metropolis}

\usepackage{pgfpages} % for notes layouts

% --- bibliography ---
\usepackage[round,authoryear]{natbib}

% --- Choose ONE of these toggles when compiling ---
% \setbeameroption{hide notes}                        % audience version (default)
% \setbeameroption{show notes}                        % notes printed under each slide
% \setbeameroption{show only notes}                   % notes-only PDF (for printing)
% \setbeameroption{show notes on second screen=right} % slide + notes side-by-side

% --- shortcuts ---
\newcommand{\E}{\mathbb{E}}
\newcommand{\R}{\mathbb{R}}
\newcommand{\btheta}{\boldsymbol{\theta}}
\newcommand{\bx}{\boldsymbol{x}}

% --- title info ---
\title{Progress meeting}
\subtitle{Closed form solutions of the NTK for 1-hidden layer ReLU networks}
\author{Shreyas Kalvankar}
\institute{TU Delft}
\date{\today}

\begin{document}

\maketitle

% -----------------------
\begin{frame}{ReLU arc-cosine kernel (covariance form)}
	Let $x,y\in\R^d$, $w\sim\mathcal N(0,I_d)$, $\sigma(t)=\max\{0,t\}$.
	\[
		K(x,y)\;:=\;\E\big[\sigma(w^\top x)\sigma(w^\top y)\big].
	\]
	Let $\theta\in[0,\pi]$ be the angle between $x$ and $y$:
	\[
		\cos\theta=\frac{\langle x,y\rangle}{\|x\|\|y\|}.
	\]
	Then the closed form is
	\[
		\boxed{
			K(x,y)=\frac{\|x\|\|y\|}{2\pi}\Big(\sin\theta+(\pi-\theta)\cos\theta\Big).
		}
	\]
	\note{I’m using the standard 2D reduction to span\{x,y\} and polar integration. I’ll keep the derivation on paper if needed.}
\end{frame}

% -----------------------
\begin{frame}{Restriction to the circle and convolution operator}
	Restrict inputs to $S^1$: $x(\phi)=(\cos\phi,\sin\phi)$, $\phi\in[0,2\pi)$.
	Then
	\[
		\langle x(\phi),x(\psi)\rangle=\cos(\phi-\psi),\qquad \delta:=\phi-\psi.
	\]
	Hence $K(\phi,\psi)=g(\delta)$ is \emph{shift-invariant} (convolution kernel).
	Define the operator
	\[
		(Tf)(\phi)=\int_0^{2\pi}K(\phi,\psi)f(\psi)\frac{d\psi}{2\pi}
		= (g*f)(\phi).
	\]
	\note{This makes the Fourier basis the right eigenbasis automatically.}
\end{frame}

% -----------------------
\begin{frame}{Fourier diagonalization}
	Because $T$ is convolution on $S^1$:
	\[
		T\cos(k\phi)=\lambda_k\cos(k\phi),\qquad
		T\sin(k\phi)=\lambda_k\sin(k\phi),
	\]
	with eigenvalues given by Fourier coefficients
	\[
		\lambda_k=\int_{0}^{2\pi} g(\delta)\cos(k\delta)\frac{d\delta}{2\pi},\qquad k\ge 0.
	\]
	\note{So the entire spectral story reduces to computing these scalar integrals.}
\end{frame}

% -----------------------
\begin{frame}{Closed-form eigenvalues: $k=1$, $k=2$, and parity}
	For the ReLU arc-cosine kernel restricted to $S^1$ (above),
	the eigenvalues admit closed forms:

	\[
		\boxed{\lambda_1=\frac18.}
	\]

	For $k>1$:
	\[
		\boxed{
			\lambda_k=
			\frac{1+(-1)^k}{\pi^2\,(k^2-1)^2}
			=
			\begin{cases}
				\displaystyle \frac{2}{\pi^2\,(k^2-1)^2}, & k \text{ even}, \\[6pt]
				0,                                        & k \text{ odd}.
			\end{cases}}
	\]

	In particular
	\[
		\lambda_2=\frac{2}{9\pi^2},\qquad
		\lambda_3=0.
	\]
	\note{Key takeaway: after $k=1$, only even modes have nonzero eigenvalues; odd modes are annihilated.}
\end{frame}

% -----------------------
\begin{frame}{Consequence: no-bias kernel cannot learn odd modes}
	Consider gradient flow for squared loss under the operator:
	\[
		\partial_t r_t = -Tr_t,\qquad r_t := f_t - f^\star.
	\]
	In Fourier:
	\[
		\partial_t \hat r_k(t)=-\lambda_k \hat r_k(t)
		\quad\Rightarrow\quad
		\hat r_k(t)=e^{-\lambda_k t}\hat r_k(0).
	\]

	Thus any component in a zero-eigenvalue eigenspace is \emph{not learned}:
	\[
		\lambda_k=0\;\Rightarrow\; \hat r_k(t)\equiv \hat r_k(0).
	\]

	For this kernel: all odd $k>1$ satisfy $\lambda_k=0$ $\;\Rightarrow\;$
	odd modes ($k=3,5,\dots$) do not decay.
	\note{This gives a clean, explicit “spectral obstruction” story without invoking finite width yet.}
\end{frame}

% -----------------------
\begin{frame}{Full NTK with learnable bias (weights + output layer + bias)}
	Use the standard 1-hidden-layer model
	\[
		f(x)=\frac{1}{\sqrt m}\sum_{i=1}^m a_i\,\sigma(w_i^\top x+b_i),
		\qquad
		a_i\sim\mathcal N(0,1),\; w_i\sim\mathcal N(0,I),\; b_i(0)=0.
	\]
	(Initialize $b_i=0$, but keep $b_i$ learnable.)

	The limiting NTK decomposes by parameter blocks:
	\[
		\Theta(x,x')=
		\E[\sigma(u)\sigma(u')]
		+\E[a^2\sigma'(u)\sigma'(u')]\,(x^\top x')
		+\E[a^2\sigma'(u)\sigma'(u')],
	\]
	where $u=w^\top x+b$, $u'=w^\top x'+b$ and $\sigma'(t)=\mathbf 1_{\{t>0\}}$.

	Since $\E[a^2]=1$:
	\[
		\boxed{\Theta(x,x')=K_0(x,x') + (x^\top x' + 1)\,K_1(x,x')}
	\]
	with
	\[
		K_0(x,x')=\E[\sigma(u)\sigma(u')],\qquad
		K_1(x,x')=\E[\mathbf 1_{\{u>0\}}\mathbf 1_{\{u'>0\}}].
	\]
	\note{The extra “$+1$” factor is exactly the bias-gradient contribution; this is the term missing in the no-bias story.}
\end{frame}

% -----------------------
\begin{frame}{Next steps}
	\begin{itemize}
		\item Restrict the \emph{full} NTK $\Theta$ to $S^1$ and compute its Fourier eigenvalues in closed form.
		\item Use those eigenvalues to obtain explicit mode-wise decay rates:
		      \[
			      \hat r_k(t)=e^{-\Lambda_k t}\hat r_k(0).
		      \]
		\item Compare parity behavior (do odd modes become learnable once bias term is included?).
	\end{itemize}
	\note{This is the key bridge to the “closed form NTK dynamics” statement and later finite-sample perturbations.}
\end{frame}
% -----------------------


% -----------------------
\end{document}
